\documentclass[UTF8,fontset=none,11pt,a4paper]{ctexart}
\usepackage{ipara-handbook}
\usepackage{amsmath,booktabs,tabularx,hyperref}
\usetikzlibrary{arrows.meta,positioning,fit,calc}

\handbooksetup{OS-B03 Process × File Reference}{408 · OS Internal Bridge}{FD · OFD · Inode · Lifetime}

\tikzset{
  fdnode/.style={draw=iparaDeep,fill=iparaSoft,rounded corners=1.5mm,align=center,minimum width=2.4cm,minimum height=0.85cm,font=\small},
  ofdnode/.style={draw=iparaGreen,fill=white,rounded corners=1.5mm,align=center,minimum width=2.6cm,minimum height=0.85cm,font=\small},
  inodenode/.style={draw=iparaDeep,double,double distance=1.2pt,fill=white,rounded corners=1.5mm,align=center,minimum width=2.6cm,minimum height=0.85cm,font=\small},
  warnnode/.style={draw=iparaAccent,fill=white,rounded corners=1.5mm,align=center,minimum width=2.4cm,minimum height=0.85cm,font=\small},
  refedge/.style={-{Latex[length=2mm]},thick,draw=iparaDeep},
  dangeredge/.style={-{Latex[length=2mm]},thick,dashed,draw=iparaAccent}
}

\begin{document}
\handbooktitle{Process × File Reference}{文件描述符不是文件：三级解耦架构怎样决定 fork、dup、close 与 unlink}{408 · OS-B03}

\begin{corebox}{Mother Interface}
\[
\boxed{
\begin{aligned}
\text{Task / Process} &\xrightarrow{\text{descriptor binding}} \text{Open Instance / OFD} \\
&\xrightarrow{\text{object reference}} \text{File Object} \\
&\xrightarrow{\text{FS mapping}} \text{Persistent Data / Metadata}
\end{aligned}
}
\]
本 Bridge 专门负责厘清“进程侧句柄绑定、系统打开实例、文件对象”之间的引用交接契约；路径解析、文件对象生命周期的最终回收协议与块分配均由 OS-06/07 独立拥有。
\end{corebox}

\begin{methodbox}{先定义接口两侧的词}
\begin{itemize}
  \item \kw{文件描述符 (fd)}：进程可见的整数句柄，用来选择本进程描述符表中的一个绑定。它不是文件身份，也不直接等于 inode。
  \item \kw{Descriptor Table / Binding（描述符表/绑定）}：进程级状态，记录某个 fd 当前引用哪个打开实例，以及少量 descriptor-local 属性。具体是数组、分层表还是其他结构属于实现。
  \item \kw{Open File Description / Open Instance（OFD/打开实例）}：一次“打开”产生的系统级运行状态，典型地拥有当前文件位置、打开状态和对底层文件对象的引用。多个 fd 可以共享同一个打开实例。
  \item \kw{File Object（文件对象）}：由文件系统标识和维护的持久对象；经典 UNIX 可用 inode 理解其身份与元数据。身份只在相应文件系统实例/命名空间契约下解释；\textbf{文件名属于命名关系，不是 inode 本体的一部分}。
\end{itemize}
\end{methodbox}

\begin{methodbox}{Owners 与职责边界}
\begin{itemize}
  \item \kw{左侧 Owner (OS-01/02 进程控制)：}拥有进程生命周期，以及 \code{fork()/exec()/exit()} 对进程资源引用的高层交接语义。
  \item \kw{右侧 Owner (OS-06/07 文件系统)：}拥有 pathname、打开实例、文件对象、命名引用、数据块映射与最终回收/持久化协议。
  \item \kw{本 Bridge Owns：}进程描述符绑定怎样跨 \code{fork/dup/close} 与打开实例交接，以及“共享同一打开实例”怎样推出共享当前文件位置等运行状态。
\end{itemize}
\end{methodbox}

\section{核心机制：引用的三级解耦数据结构}

稳定模型不是记 Linux 结构体名字，而是保留三层独立身份：
\begin{enumerate}
  \item \textbf{进程句柄层}：当前进程用 fd 选择一个描述符绑定；不同进程可以使用相同整数 fd 指向完全不同的对象；
  \item \textbf{打开实例层}：一次打开产生运行时状态。若多个描述符绑定引用同一 OFD，它们共享该 OFD 所拥有的当前文件位置等状态；
  \item \textbf{文件对象层}：多个独立 OFD 仍可以引用同一个文件对象，但各自保持独立的打开实例状态。
\end{enumerate}
这三层使“共享文件对象”与“共享当前 offset”成为两个不同问题。

\section{母模型：独立打开、\texttt{fork()} 继承与 \texttt{dup()} 拓扑}

\begin{center}
\begin{tikzpicture}[x=3.40cm,y=1.55cm]
  % Level 1: fd
  \node[fdnode] (pfd) at (0, 0.8) {Parent\\[-1pt]\scriptsize fd = 3};
  \node[fdnode] (cfd) at (0, -0.8) {Child (fork)\\[-1pt]\scriptsize fd = 3};
  \node[fdnode] (dupfd) at (0, 0) {Parent (dup)\\[-1pt]\scriptsize fd = 4};
  \node[warnnode] (otherfd) at (0, -1.8) {Other Proc\\[-1pt]\scriptsize fd = 7};

  % Level 2: OFD
  \node[ofdnode] (ofd1) at (1.4, 0.4) {OFD \#1\\[-1pt]\scriptsize \code{f\_pos=100, f\_count=3}};
  \node[ofdnode] (ofd2) at (1.4, -1.8) {OFD \#2\\[-1pt]\scriptsize \code{f\_pos=0, f\_count=1}};

  % Level 3: Inode
  \node[inodenode] (inode) at (2.8, -0.7) {Inode Object\\[-1pt]\scriptsize inode=8848 (\code{nlink=1})};

  % Connections
  \draw[refedge] (pfd) -- (ofd1);
  \draw[refedge] (dupfd) -- (ofd1);
  \draw[refedge] (cfd) -- (ofd1);
  \draw[refedge] (otherfd) -- (ofd2);

  \draw[refedge] (ofd1) -- (inode);
  \draw[refedge] (ofd2) -- (inode);
\end{tikzpicture}
\end{center}

\subsection{读写位置（\texttt{f\_pos}）共享判定黄金法则}
\begin{itemize}
  \item \kw{共享同一个 OFD（游标共同推进）}：
  \begin{itemize}
    \item \code{fork()}：子进程复制父进程的 \code{fd\_table}，父子相同 $fd$ 指向同一 OFD，父进程读写会影响子进程的后续读取位置；
    \item \code{dup(oldfd) / dup2()}：同一进程产生新 $fd$，指向同一 OFD，两个 $fd$ 操作同一个读写游标。
  \end{itemize}
  \item \kw{拥有各自独立的 OFD（游标互不影响）}：
  \begin{itemize}
    \item 两个进程（或同一进程两次）独立调用 \code{open("same.txt", ...)}，各自生成独立 OFD，拥有独立的 \code{f\_pos}，但最终指向同一个 Inode。
  \end{itemize}
\end{itemize}

\section{常见操作怎样改变三层关系}

\begin{center}
\small
\renewcommand{\arraystretch}{1.25}
\begin{tabularx}{\linewidth}{>{\bfseries}p{0.18\linewidth} X X X}
\toprule
操作 API & 进程描述符绑定 & 打开实例 / OFD & 文件对象 / 命名关系 \\
\midrule
\code{open()} & 建立一个新的 fd 绑定 & 通常建立一个新的打开实例，拥有自己的 current offset 等状态 & pathname 被解析到一个文件对象；创建/截断等效果按具体 flags 与 FS 契约处理 \\
\code{fork()} & 子进程得到描述符绑定的语义副本 & 父子对应绑定继续引用同一打开实例，因此共享其 current offset & 不创造新的文件对象或硬链接 \\
\code{dup/dup2()} & 在同一进程建立另一个 fd 绑定 & 新旧 fd 引用同一打开实例 & 文件对象不因 dup 被复制 \\
\code{close(fd)} & 删除当前进程的一条 fd 绑定 & 该打开实例少一个引用；是否最终销毁由剩余引用决定 & 不等于删除命名关系；文件对象回收由 OS-06/07 生命周期协议决定 \\
\code{link()} & 与既有 fd 绑定无直接关系 & 不要求创建新打开实例 & 增加一个指向同一文件对象的命名引用 \\
\code{unlink()} & 既有 fd 绑定仍可继续存在 & 既有打开实例仍可引用该对象 & 删除一条命名引用；对象何时最终可回收取决于命名引用与打开引用等生命周期条件 \\
\bottomrule
\end{tabularx}
\end{center}

\section{经典反直觉案例：\texttt{unlink()} 正在被打开的文件}

若进程 P 已经打开 \code{/tmp/temp.dat}，此时另一执行路径执行 \code{unlink("/tmp/temp.dat")}：
\begin{enumerate}
  \item 文件系统删除的是该 pathname 对文件对象的一条\kw{命名引用}；若没有其他同名/硬链接，新 pathname lookup 将无法再由这条名字找到对象；
  \item P 的 fd 绑定和打开实例并不会因为名字消失而自动失效；
  \item 因此 P 仍可通过既有 fd 继续按该打开实例的权限与 offset 访问对象；
  \item 只有当文件系统判断\kw{命名引用和打开引用等所有维持对象生命的关系都已消失}后，该对象才进入最终删除/空间回收流程；具体延迟回收、日志提交和缓存写回由 OS-06/07 Own。
\end{enumerate}

\begin{warnbox}{Anti-Bridge：三个必须阻断的错误认知}
\begin{itemize}
  \item \(\text{fd} \neq \text{Inode}\)：$fd$ 只是进程私有整数代号；Inode 才是持久文件对象。
  \item \(\text{close(fd)} \neq \text{Delete File}\)：\code{close()} 仅仅解除了当前进程的一个引用句柄，文件本身依然完好保存在文件系统中。
  \item \(\text{unlink()} \neq \text{Immediate Disk Reclamation}\)：删除 pathname 只改变命名关系；对象何时成为可回收对象，要继续检查是否还有其他命名引用、打开引用以及文件系统自身尚未完成的持久化/回收状态。
\end{itemize}
\end{warnbox}

\begin{corebox}{30 秒极速调用协议}
做题遇到“fork/dup/open 多次读写、unlink 空间释放”时，严格按三层追踪：
\[
\boxed{fd \text{ 绑定到哪个打开实例？}} \xrightarrow{\text{是否共享同一 OFD/current offset}} \boxed{\text{命名引用与打开引用分别怎样变化？}}
\]
\end{corebox}

\end{document}
